% Problem record op_a40ad449c54093d7. This TeX file is derived from
% database/problems_json/op_a40ad449c54093d7.json by scripts/sync-tex.mjs; edit the JSON record
% and rerun the script rather than editing this file.
\section{Optimal random-unitary decomposition of symmetric Werner--Holevo channels}
\label{sec:op_a40ad449c54093d7}

\paragraph{Problem.}
For every odd integer $d\geq5$, do there exist $r=d(d+1)/2$ unitary
operators $U_1,\ldots,U_r$ on $\mathbb C^d$ such that
\begin{equation}
  \Phi_d(X):=\frac{\operatorname{Tr}(X)I_d+X^{\mathsf T}}{d+1}
  =\frac1r\sum_{j=1}^{r}U_j XU_j^\dagger
  \qquad\text{for every }X\in\mathcal L(\mathbb C^d)?
  \label{eq:a40ad-wh-decomposition}
\end{equation}
The transpose in Eq.~\eqref{eq:a40ad-wh-decomposition} is taken in a fixed
orthonormal basis.

\subsection*{Status}

Unsolved

\subsection*{Source}

Girard et al. explicitly conjecture that the symmetric Werner--Holevo
channel has mixed-unitary rank equal to its Choi rank; see Section~6,
final paragraph after Theorem~23, p.~29 of the preprint
\sourcecite{ref:a40ad-girard}{GLL+22}.

\subsection*{Progress}

\begin{itemize}
  \item Let $N(\Phi_d)$ denote the minimum number of unitary conjugations in a
convex decomposition. Its Choi rank gives $N(\Phi_d)\geq r$; for odd $d$,
Theorem~22 gives $N(\Phi_d)\leq d(d+3)/2$
\sourcecite{ref:a40ad-girard}{GLL+22}.

  \item The optimum $N(\Phi_d)=r$ holds for every even $d$ and for $d=3$
(Theorems~21 and~23) \sourcecite{ref:a40ad-girard}{GLL+22}.
Levick and Rahaman prove it for every prime $d\equiv7\pmod8$
(Theorem~3.8 and Corollary~3.12) \sourcecite{ref:a40ad-levick}{LR21}.

  \item Theorem~19 identifies $N(\Phi_d)=r$ with the existence of a
Hilbert--Schmidt orthogonal basis of the complex symmetric matrices
consisting of unitaries. Every optimal $r$-term decomposition necessarily
has weights $1/r$, so the uniform formulation in
Eq.~\eqref{eq:a40ad-wh-decomposition} is equivalent to the rank conjecture
\sourcecite{ref:a40ad-girard}{GLL+22}.
\end{itemize}

\subsection*{References}

\begin{enumerate}[label={},leftmargin=4.8em,labelsep=0.5em]
  \footnotesize
  \item[\textup{[GLL+22]}]\label{ref:a40ad-girard}
  M. Girard, D. Leung, J. Levick, C.-K. Li, V. Paulsen, Y. T. Poon,
and J. Watrous, ``On the Mixed-Unitary Rank of Quantum Channels,''
\emph{Communications in Mathematical Physics} \textbf{394}, 919--951 (2022).
\href{https://doi.org/10.1007/s00220-022-04412-y}{doi:10.1007/s00220-022-04412-y};
\href{https://arxiv.org/abs/2003.14405}{arXiv:2003.14405}.

  \item[\textup{[LR21]}]\label{ref:a40ad-levick}
  J. Levick and M. Rahaman, ``An extension of Bravyi--Smolin's
construction for UMEBs,'' \emph{Quantum Information Processing}
\textbf{20}, 369 (2021).
\href{https://doi.org/10.1007/s11128-021-03312-9}{doi:10.1007/s11128-021-03312-9};
\href{https://arxiv.org/abs/2105.00975v2}{arXiv:2105.00975v2}.
\end{enumerate}

\subsection*{Comment}

The unresolved target is the universal odd-dimensional optimum.
Numerical decompositions do not supply an exact proof. Literature audit:
9 September 2026. Searches under the channel, rank, and symmetric-unitary-basis
formulations found no full resolution;
indexed literature searches cannot exclude unindexed work. The cited partial
results are peer-reviewed.

\subsection*{Field}

Quantum Communication

\subsection*{Topic}

Quantum channel structure

\subsection*{ID}

\texttt{op\_a40ad449c54093d7}
